\documentclass[11pt,a4paper]{article}

\newcommand{\doctitle}{Gauge Historical Daily Synthetic Generation Model (MKM-GHD-001)}
\newcommand{\docsubtitle}{Wolfram Oracle Benchmark: the Series' Statistics Computed Exactly}
\newcommand{\docversion}{0.1}
\newcommand{\docdate}{28-October-2026}
\newcommand{\docauthor}{David K Kelly}

% MKM Research Labs — Shared Document Header
% Usage: % MKM Research Labs — Shared Document Header
% Usage: % MKM Research Labs — Shared Document Header
% Usage: \input{../shared/mkm_header}  (before \begin{document})
%
% Required commands before \input:
%   \newcommand{\doctitle}{...}
%   \newcommand{\docsubtitle}{...}
%   \newcommand{\docversion}{1.0}
%   \newcommand{\docdate}{13-February-2026}
%   \newcommand{\docauthor}{David K Kelly}

% ─── Packages ───────────────────────────────────────────────────────────────
\usepackage[utf8]{inputenc}
\usepackage{amsmath,amsfonts,amssymb}
\usepackage{listings}
\usepackage{color,xcolor}
\usepackage{graphicx}
\usepackage{hyperref}
\usepackage[toc,page]{appendix}
\usepackage{tabto}
\usepackage{float}
\usepackage{comment}
\usepackage[en-GB]{datetime2}
\usepackage{setspace}
\usepackage{subcaption}
\usepackage{booktabs}
\usepackage{longtable}
\usepackage{adjustbox}
\usepackage{enumitem}
\usepackage{fancyhdr}
\usepackage{lastpage}

% ─── Page geometry ──────────────────────────────────────────────────────────
\onehalfspacing
\usepackage{geometry}
\geometry{a4paper, total={170mm,247mm}, left=20mm, top=25mm, bottom=25mm}

% ─── Headers and footers ───────────────────────────────────────────────────
\pagestyle{fancy}
\fancyhf{}
\fancyhead[L]{\small\textcolor{gray}{MKM Research Labs}}
\fancyhead[R]{\small\textcolor{gray}{\doctitle{} --- v\docversion}}
\fancyfoot[C]{\small\textcolor{gray}{CONFIDENTIAL}}
\fancyfoot[R]{\small\textcolor{gray}{Page \thepage\ of \pageref{LastPage}}}
\renewcommand{\headrulewidth}{0.4pt}
\renewcommand{\footrulewidth}{0.4pt}

% ─── Colours ────────────────────────────────────────────────────────────────
\definecolor{mkmblue}{HTML}{1565C0}
\definecolor{codegreen}{rgb}{0,0.6,0}
\definecolor{codegray}{rgb}{0.5,0.5,0.5}
\definecolor{codepurple}{rgb}{0.58,0,0.82}
\definecolor{backcolour}{rgb}{0.95,0.95,0.92}

% ─── Hyperref setup ─────────────────────────────────────────────────────────
\hypersetup{
    colorlinks=true,
    linkcolor=mkmblue,
    citecolor=mkmblue,
    urlcolor=mkmblue
}

% ─── Code listings ──────────────────────────────────────────────────────────
\lstdefinestyle{mystyle}{
    backgroundcolor=\color{backcolour},
    commentstyle=\color{codegreen},
    keywordstyle=\color{magenta},
    numberstyle=\tiny\color{codegray},
    stringstyle=\color{codepurple},
    basicstyle=\ttfamily\footnotesize,
    breakatwhitespace=false,
    breaklines=true,
    captionpos=b,
    keepspaces=true,
    numbers=left,
    numbersep=5pt,
    showspaces=false,
    showstringspaces=false,
    showtabs=false,
    tabsize=2
}
\lstset{style=mystyle}

% ─── Section formatting ────────────────────────────────────────────────────
\setcounter{tocdepth}{3}
\setcounter{secnumdepth}{3}

% ─── Shared macros ──────────────────────────────────────────────────────────
\newcommand{\mkmtitlepage}{%
    \begin{titlepage}
    \begin{center}
    \vspace*{1cm}
    \begin{figure}[H]\centering
        \includegraphics[width=0.4\textwidth]{../shared/Logos/MKM_logo.png}
    \end{figure}
    \vspace{1.5cm}
    {\Huge\bfseries\textcolor{mkmblue}{\doctitle}\par}
    \vspace{0.5cm}
    {\Large\textcolor{gray}{\docsubtitle}\par}
    \vspace{2cm}
    {\large MKM Research Labs\par}
    \vspace{0.5cm}
    {\large Version \docversion{} --- \docdate\par}
    \vspace{0.5cm}
    {\large \docauthor\par}
    \vspace{2cm}
    \textcolor{red}{\textbf{CONFIDENTIAL --- PROPRIETARY}}
    \end{center}
    \end{titlepage}
}

\newcommand{\mkmlegalpage}{%
    \clearpage
    \thispagestyle{empty}
    \begin{center}
    \large\textbf{Legal Notice}
    \vspace{1em}

    \textbf{PROPRIETARY AND CONFIDENTIAL}
    \end{center}

    \noindent This document, including all algorithms, models, methodology, formulas,
    calculations, and intellectual concepts contained herein, constitutes the exclusive
    intellectual property and confidential trade secrets of MKM Research Labs (``MKM'').

    \vspace{0.5em}
    \noindent Any usage, reproduction, distribution, modification, or disclosure of this
    document or any portion thereof, without the express prior written authorisation from
    MKM Research Labs is strictly prohibited and constitutes an infringement of intellectual
    property rights.

    \vspace{0.5em}
    \noindent The document is provided on an ``as is'' basis. MKM Research Labs makes no
    representations or warranties of any kind, express or implied, regarding its accuracy,
    reliability, or suitability for any particular purpose.

    \vspace{1em}
    \begin{center}
    All rights reserved. \copyright\ 2021--2026 MKM Research Labs.

    \vspace{0.5em}
    Governed by the laws of the United Kingdom.
    \end{center}
    \clearpage
}

\newcommand{\mkmhistorypage}[1]{%
    \section{Document History}
    \begin{table}[ht]
        \centering
        \begin{tabular}{l|l|l|l}
            \textbf{Date} & \textbf{Version} & \textbf{Description} & \textbf{Author} \\
            \hline
            #1
        \end{tabular}
        \label{tab:revision_history}
    \end{table}
    \clearpage
}

% ─── Bank validation protocol section commands ──────────────────────────────
% Use these to ensure every model doc has the required sections.
% They are just styled \section calls — content is provided by each doc.

\newcommand{\mkmexecsummary}{\section{Executive Summary}}
\newcommand{\mkmpurpose}{\section{Model Purpose and Scope}}
\newcommand{\mkmframework}{\section{Mathematical Framework}}
\newcommand{\mkminputs}{\section{Input Parameters and Calibration}}
\newcommand{\mkmimplementation}{\section{Implementation Details}}
\newcommand{\mkmsensitivity}{\section{Sensitivity Analysis}}
\newcommand{\mkmvalidation}{\section{Validation and Backtesting}}
\newcommand{\mkmtestresults}{\subsection{Automated Test Results}\input{test_results}}
\newcommand{\mkmlimitations}{\section{Model Limitations and Known Weaknesses}}

\newcommand{\mkmgovernance}{\section{Governance Validation Questions}
\noindent The following nine questions are mandated by Chapter~5 of the
\emph{Model Risk Governance for Vendors} handbook and must be explicitly
addressed for every model in the inventory. Response status is tracked in the
Model Governance Dashboard (MRC portal). \par\medskip}

\newcommand{\mkmvqitem}[3]{\subsection{Q#1 --- #2}
\noindent\textit{``#3''}\par\medskip\noindent}

\newcommand{\mkmchangelog}{\section{Change History}}
  (before \begin{document})
%
% Required commands before \input:
%   \newcommand{\doctitle}{...}
%   \newcommand{\docsubtitle}{...}
%   \newcommand{\docversion}{1.0}
%   \newcommand{\docdate}{13-February-2026}
%   \newcommand{\docauthor}{David K Kelly}

% ─── Packages ───────────────────────────────────────────────────────────────
\usepackage[utf8]{inputenc}
\usepackage{amsmath,amsfonts,amssymb}
\usepackage{listings}
\usepackage{color,xcolor}
\usepackage{graphicx}
\usepackage{hyperref}
\usepackage[toc,page]{appendix}
\usepackage{tabto}
\usepackage{float}
\usepackage{comment}
\usepackage[en-GB]{datetime2}
\usepackage{setspace}
\usepackage{subcaption}
\usepackage{booktabs}
\usepackage{longtable}
\usepackage{adjustbox}
\usepackage{enumitem}
\usepackage{fancyhdr}
\usepackage{lastpage}

% ─── Page geometry ──────────────────────────────────────────────────────────
\onehalfspacing
\usepackage{geometry}
\geometry{a4paper, total={170mm,247mm}, left=20mm, top=25mm, bottom=25mm}

% ─── Headers and footers ───────────────────────────────────────────────────
\pagestyle{fancy}
\fancyhf{}
\fancyhead[L]{\small\textcolor{gray}{MKM Research Labs}}
\fancyhead[R]{\small\textcolor{gray}{\doctitle{} --- v\docversion}}
\fancyfoot[C]{\small\textcolor{gray}{CONFIDENTIAL}}
\fancyfoot[R]{\small\textcolor{gray}{Page \thepage\ of \pageref{LastPage}}}
\renewcommand{\headrulewidth}{0.4pt}
\renewcommand{\footrulewidth}{0.4pt}

% ─── Colours ────────────────────────────────────────────────────────────────
\definecolor{mkmblue}{HTML}{1565C0}
\definecolor{codegreen}{rgb}{0,0.6,0}
\definecolor{codegray}{rgb}{0.5,0.5,0.5}
\definecolor{codepurple}{rgb}{0.58,0,0.82}
\definecolor{backcolour}{rgb}{0.95,0.95,0.92}

% ─── Hyperref setup ─────────────────────────────────────────────────────────
\hypersetup{
    colorlinks=true,
    linkcolor=mkmblue,
    citecolor=mkmblue,
    urlcolor=mkmblue
}

% ─── Code listings ──────────────────────────────────────────────────────────
\lstdefinestyle{mystyle}{
    backgroundcolor=\color{backcolour},
    commentstyle=\color{codegreen},
    keywordstyle=\color{magenta},
    numberstyle=\tiny\color{codegray},
    stringstyle=\color{codepurple},
    basicstyle=\ttfamily\footnotesize,
    breakatwhitespace=false,
    breaklines=true,
    captionpos=b,
    keepspaces=true,
    numbers=left,
    numbersep=5pt,
    showspaces=false,
    showstringspaces=false,
    showtabs=false,
    tabsize=2
}
\lstset{style=mystyle}

% ─── Section formatting ────────────────────────────────────────────────────
\setcounter{tocdepth}{3}
\setcounter{secnumdepth}{3}

% ─── Shared macros ──────────────────────────────────────────────────────────
\newcommand{\mkmtitlepage}{%
    \begin{titlepage}
    \begin{center}
    \vspace*{1cm}
    \begin{figure}[H]\centering
        \includegraphics[width=0.4\textwidth]{../shared/Logos/MKM_logo.png}
    \end{figure}
    \vspace{1.5cm}
    {\Huge\bfseries\textcolor{mkmblue}{\doctitle}\par}
    \vspace{0.5cm}
    {\Large\textcolor{gray}{\docsubtitle}\par}
    \vspace{2cm}
    {\large MKM Research Labs\par}
    \vspace{0.5cm}
    {\large Version \docversion{} --- \docdate\par}
    \vspace{0.5cm}
    {\large \docauthor\par}
    \vspace{2cm}
    \textcolor{red}{\textbf{CONFIDENTIAL --- PROPRIETARY}}
    \end{center}
    \end{titlepage}
}

\newcommand{\mkmlegalpage}{%
    \clearpage
    \thispagestyle{empty}
    \begin{center}
    \large\textbf{Legal Notice}
    \vspace{1em}

    \textbf{PROPRIETARY AND CONFIDENTIAL}
    \end{center}

    \noindent This document, including all algorithms, models, methodology, formulas,
    calculations, and intellectual concepts contained herein, constitutes the exclusive
    intellectual property and confidential trade secrets of MKM Research Labs (``MKM'').

    \vspace{0.5em}
    \noindent Any usage, reproduction, distribution, modification, or disclosure of this
    document or any portion thereof, without the express prior written authorisation from
    MKM Research Labs is strictly prohibited and constitutes an infringement of intellectual
    property rights.

    \vspace{0.5em}
    \noindent The document is provided on an ``as is'' basis. MKM Research Labs makes no
    representations or warranties of any kind, express or implied, regarding its accuracy,
    reliability, or suitability for any particular purpose.

    \vspace{1em}
    \begin{center}
    All rights reserved. \copyright\ 2021--2026 MKM Research Labs.

    \vspace{0.5em}
    Governed by the laws of the United Kingdom.
    \end{center}
    \clearpage
}

\newcommand{\mkmhistorypage}[1]{%
    \section{Document History}
    \begin{table}[ht]
        \centering
        \begin{tabular}{l|l|l|l}
            \textbf{Date} & \textbf{Version} & \textbf{Description} & \textbf{Author} \\
            \hline
            #1
        \end{tabular}
        \label{tab:revision_history}
    \end{table}
    \clearpage
}

% ─── Bank validation protocol section commands ──────────────────────────────
% Use these to ensure every model doc has the required sections.
% They are just styled \section calls — content is provided by each doc.

\newcommand{\mkmexecsummary}{\section{Executive Summary}}
\newcommand{\mkmpurpose}{\section{Model Purpose and Scope}}
\newcommand{\mkmframework}{\section{Mathematical Framework}}
\newcommand{\mkminputs}{\section{Input Parameters and Calibration}}
\newcommand{\mkmimplementation}{\section{Implementation Details}}
\newcommand{\mkmsensitivity}{\section{Sensitivity Analysis}}
\newcommand{\mkmvalidation}{\section{Validation and Backtesting}}
\newcommand{\mkmtestresults}{\subsection{Automated Test Results}\input{test_results}}
\newcommand{\mkmlimitations}{\section{Model Limitations and Known Weaknesses}}

\newcommand{\mkmgovernance}{\section{Governance Validation Questions}
\noindent The following nine questions are mandated by Chapter~5 of the
\emph{Model Risk Governance for Vendors} handbook and must be explicitly
addressed for every model in the inventory. Response status is tracked in the
Model Governance Dashboard (MRC portal). \par\medskip}

\newcommand{\mkmvqitem}[3]{\subsection{Q#1 --- #2}
\noindent\textit{``#3''}\par\medskip\noindent}

\newcommand{\mkmchangelog}{\section{Change History}}
  (before \begin{document})
%
% Required commands before \input:
%   \newcommand{\doctitle}{...}
%   \newcommand{\docsubtitle}{...}
%   \newcommand{\docversion}{1.0}
%   \newcommand{\docdate}{13-February-2026}
%   \newcommand{\docauthor}{David K Kelly}

% ─── Packages ───────────────────────────────────────────────────────────────
\usepackage[utf8]{inputenc}
\usepackage{amsmath,amsfonts,amssymb}
\usepackage{listings}
\usepackage{color,xcolor}
\usepackage{graphicx}
\usepackage{hyperref}
\usepackage[toc,page]{appendix}
\usepackage{tabto}
\usepackage{float}
\usepackage{comment}
\usepackage[en-GB]{datetime2}
\usepackage{setspace}
\usepackage{subcaption}
\usepackage{booktabs}
\usepackage{longtable}
\usepackage{adjustbox}
\usepackage{enumitem}
\usepackage{fancyhdr}
\usepackage{lastpage}

% ─── Page geometry ──────────────────────────────────────────────────────────
\onehalfspacing
\usepackage{geometry}
\geometry{a4paper, total={170mm,247mm}, left=20mm, top=25mm, bottom=25mm}

% ─── Headers and footers ───────────────────────────────────────────────────
\pagestyle{fancy}
\fancyhf{}
\fancyhead[L]{\small\textcolor{gray}{MKM Research Labs}}
\fancyhead[R]{\small\textcolor{gray}{\doctitle{} --- v\docversion}}
\fancyfoot[C]{\small\textcolor{gray}{CONFIDENTIAL}}
\fancyfoot[R]{\small\textcolor{gray}{Page \thepage\ of \pageref{LastPage}}}
\renewcommand{\headrulewidth}{0.4pt}
\renewcommand{\footrulewidth}{0.4pt}

% ─── Colours ────────────────────────────────────────────────────────────────
\definecolor{mkmblue}{HTML}{1565C0}
\definecolor{codegreen}{rgb}{0,0.6,0}
\definecolor{codegray}{rgb}{0.5,0.5,0.5}
\definecolor{codepurple}{rgb}{0.58,0,0.82}
\definecolor{backcolour}{rgb}{0.95,0.95,0.92}

% ─── Hyperref setup ─────────────────────────────────────────────────────────
\hypersetup{
    colorlinks=true,
    linkcolor=mkmblue,
    citecolor=mkmblue,
    urlcolor=mkmblue
}

% ─── Code listings ──────────────────────────────────────────────────────────
\lstdefinestyle{mystyle}{
    backgroundcolor=\color{backcolour},
    commentstyle=\color{codegreen},
    keywordstyle=\color{magenta},
    numberstyle=\tiny\color{codegray},
    stringstyle=\color{codepurple},
    basicstyle=\ttfamily\footnotesize,
    breakatwhitespace=false,
    breaklines=true,
    captionpos=b,
    keepspaces=true,
    numbers=left,
    numbersep=5pt,
    showspaces=false,
    showstringspaces=false,
    showtabs=false,
    tabsize=2
}
\lstset{style=mystyle}

% ─── Section formatting ────────────────────────────────────────────────────
\setcounter{tocdepth}{3}
\setcounter{secnumdepth}{3}

% ─── Shared macros ──────────────────────────────────────────────────────────
\newcommand{\mkmtitlepage}{%
    \begin{titlepage}
    \begin{center}
    \vspace*{1cm}
    \begin{figure}[H]\centering
        \includegraphics[width=0.4\textwidth]{../shared/Logos/MKM_logo.png}
    \end{figure}
    \vspace{1.5cm}
    {\Huge\bfseries\textcolor{mkmblue}{\doctitle}\par}
    \vspace{0.5cm}
    {\Large\textcolor{gray}{\docsubtitle}\par}
    \vspace{2cm}
    {\large MKM Research Labs\par}
    \vspace{0.5cm}
    {\large Version \docversion{} --- \docdate\par}
    \vspace{0.5cm}
    {\large \docauthor\par}
    \vspace{2cm}
    \textcolor{red}{\textbf{CONFIDENTIAL --- PROPRIETARY}}
    \end{center}
    \end{titlepage}
}

\newcommand{\mkmlegalpage}{%
    \clearpage
    \thispagestyle{empty}
    \begin{center}
    \large\textbf{Legal Notice}
    \vspace{1em}

    \textbf{PROPRIETARY AND CONFIDENTIAL}
    \end{center}

    \noindent This document, including all algorithms, models, methodology, formulas,
    calculations, and intellectual concepts contained herein, constitutes the exclusive
    intellectual property and confidential trade secrets of MKM Research Labs (``MKM'').

    \vspace{0.5em}
    \noindent Any usage, reproduction, distribution, modification, or disclosure of this
    document or any portion thereof, without the express prior written authorisation from
    MKM Research Labs is strictly prohibited and constitutes an infringement of intellectual
    property rights.

    \vspace{0.5em}
    \noindent The document is provided on an ``as is'' basis. MKM Research Labs makes no
    representations or warranties of any kind, express or implied, regarding its accuracy,
    reliability, or suitability for any particular purpose.

    \vspace{1em}
    \begin{center}
    All rights reserved. \copyright\ 2021--2026 MKM Research Labs.

    \vspace{0.5em}
    Governed by the laws of the United Kingdom.
    \end{center}
    \clearpage
}

\newcommand{\mkmhistorypage}[1]{%
    \section{Document History}
    \begin{table}[ht]
        \centering
        \begin{tabular}{l|l|l|l}
            \textbf{Date} & \textbf{Version} & \textbf{Description} & \textbf{Author} \\
            \hline
            #1
        \end{tabular}
        \label{tab:revision_history}
    \end{table}
    \clearpage
}

% ─── Bank validation protocol section commands ──────────────────────────────
% Use these to ensure every model doc has the required sections.
% They are just styled \section calls — content is provided by each doc.

\newcommand{\mkmexecsummary}{\section{Executive Summary}}
\newcommand{\mkmpurpose}{\section{Model Purpose and Scope}}
\newcommand{\mkmframework}{\section{Mathematical Framework}}
\newcommand{\mkminputs}{\section{Input Parameters and Calibration}}
\newcommand{\mkmimplementation}{\section{Implementation Details}}
\newcommand{\mkmsensitivity}{\section{Sensitivity Analysis}}
\newcommand{\mkmvalidation}{\section{Validation and Backtesting}}
\newcommand{\mkmtestresults}{\subsection{Automated Test Results}\input{test_results}}
\newcommand{\mkmlimitations}{\section{Model Limitations and Known Weaknesses}}

\newcommand{\mkmgovernance}{\section{Governance Validation Questions}
\noindent The following nine questions are mandated by Chapter~5 of the
\emph{Model Risk Governance for Vendors} handbook and must be explicitly
addressed for every model in the inventory. Response status is tracked in the
Model Governance Dashboard (MRC portal). \par\medskip}

\newcommand{\mkmvqitem}[3]{\subsection{Q#1 --- #2}
\noindent\textit{``#3''}\par\medskip\noindent}

\newcommand{\mkmchangelog}{\section{Change History}}

% T1 so the underscores in code identifiers print as underscores.
\usepackage[T1]{fontenc}
\usepackage{lmodern}
% Generated by tools.wolfram_oracle -- do not edit.
\newcommand{\bmGenerated}{2026-10-28}
\newcommand{\bmCodeCommit}{\texttt{7bbb63078a59}}
\newcommand{\bmCodeDirty}{}
\newcommand{\bmNumerics}{numpy 1.26.4, scipy 1.15.1, mpmath 1.4.1}
\newcommand{\bmWolfram}{15.0.0 for Mac OS X ARM (64-bit) (May 26, 2026)}
\newcommand{\bmWolframVia}{ via live kernel session (wolframclient)}
\newcommand{\bmWolframConfirms}{Wolfram's exact integrals agree with the reference's double-precision closed forms to $4.9\times10^{-10}$.}
\newcommand{\bmChecks}{7}
\newcommand{\bmOutcome}{all pass}
\newcommand{\bmReps}{200}
\newcommand{\bmWorstZ}{2.03}
\newcommand{\bmDays}{18{,}251}
\newcommand{\bmLagExact}{0.981}
\newcommand{\bmLagModel}{0.967}
\newcommand{\bmAlertZero}{16.3}
\newcommand{\bmAlertEight}{129.7}
\newcommand{\bmPerEvent}{2.7}
\newcommand{\bmSeedOne}{4181464545}
\newcommand{\bmSeedTwo}{3341949591}
\newcommand{\bmJump}{9.8}

% The shared history table sets every column to its natural width; wrap the description.
\renewcommand{\mkmhistorypage}[1]{%
    \section{Document History}
    \begin{table}[ht]
        \centering
        \begin{tabular}{l|l|p{9cm}|l}
            \textbf{Date} & \textbf{Version} & \textbf{Description} & \textbf{Author} \\
            \hline
            #1
        \end{tabular}
    \end{table}
}

\begin{document}
% Long code paths and identifiers cannot hyphenate; let the spacing stretch instead.
\sloppy

\mkmtitlepage
\mkmlegalpage

\tableofcontents
\clearpage

% =========================================================================
\mkmexecsummary
% =========================================================================

This document benchmarks the Gauge Historical Daily Synthetic Generation Model
(MKM-GHD-001), which generates a 50-year daily water-level history for each gauge
from its flood thresholds, tidal regime and recorded count of severe exceedances.
Those histories feed the gauge statistics and, through their monthly means, the
stress pipeline. It asks whether the code generates what its model document
(\texttt{gaugehd\_synthetic.pdf}, v1.0) says, and whether what the document says
about the result holds.

\paragraph{The code generates its documented process.} The model is random, so it is
judged in distribution. For six gauges (each tidal regime, with no recorded severe
exceedance and with eight), the expected number of days at or above each flood
threshold, the mean level, and the January and July departures from it are computed
exactly from the document's equations (1)--(14). Each is then compared with the mean
over \bmReps{} generated series of \bmDays{} days. Every statistic agrees within
sampling error, worst $|z|=\bmWorstZ$ (Tables~\ref{tab:bm-checks}
and~\ref{tab:bm-gauges}). \bmWolframConfirms

\paragraph{But several of the document's own statements about the output do not
hold} (Table~\ref{tab:bm-claims}):
\begin{enumerate}[leftmargin=*]
\item \textbf{Severe frequency is understated tenfold (\S4.5).}
  \texttt{FrequencyExceedLevel3} counts severe exceedances over five years (Table~2),
  but Eq.~(9) spreads that count over the whole 50-year run. A gauge with eight severe
  exceedances in five years is given eight in fifty. This is the model's main control
  on the upper tail (\S6.4).
\item \textbf{Regeneration is not reproducible (\S5.4, \S7.3).} The seed is Python's
  \texttt{hash()} of the gauge id, which is salted per process: the same id seeds
  \bmSeedOne{} in one process and \bmSeedTwo{} in another. The series also ends on
  the day it is run, so even a fixed seed gives a different calendar each day.
\item \textbf{The flood-frequency diagnostic cannot pass (\S7.1(3)).} It expects about
  $6\max(n_3,1)$ days above FloodAlert. The exact expectation is \bmAlertZero{} and
  \bmAlertEight{} for $n_3=0$ and 8, because a flood lifts the AR(1) state and the
  level stays above the alert for \bmPerEvent{} days on average (non-tidal). For the
  same reason, \S8.4's single-day spikes are not what the model produces.
\item \textbf{The lag-1 diagnostic cannot pass (\S7.1(1)).} The AR(1) noise has lag-1
  autocorrelation 0.85, but the series is dominated by its seasonal cycle. Its lag-1
  autocorrelation is \bmLagExact{} without floods, and \bmLagModel{} as generated.
\item \textbf{The spring--neap cycle restarts each year (\S3.5).} Eq.~(6) takes the
  phase from the day of the year, so it jumps \bmJump{} days each 1 January.
\end{enumerate}
The seasonal means (\S5.6), the lag-7 autocorrelations (\S6.1) and the mean levels
(\S7.2) are as stated.

\paragraph{Evidence.} \bmChecks{} checks, outcome: \bmOutcome{}
(Table~\ref{tab:bm-checks}).
\begin{itemize}[leftmargin=*,nosep]
\item Wolfram: \bmWolfram{}\bmWolframVia.
\item Code under test, at commit \bmCodeCommit{}\bmCodeDirty.
\item Numerics: \bmNumerics.
\item Generated \bmGenerated{} by \texttt{python -m tools.wolfram\_oracle ghd001}.
\end{itemize}

% =========================================================================
\mkmpurpose
% =========================================================================

\subsection{What is benchmarked}
\texttt{generate\_synthetic\_timeseries} in \texttt{src/models/statistics/synthetic.py},
imported as production imports it and called with a gauge record as
\texttt{src/port/src/gauge/gaugehd/synthetic.py} builds one, on the document's default
thresholds (3, 4 and 5\,m) and a historical high of 5.5\,m.

\subsection{What is not}
The historical-high injection (\S3.8) and the 0.1\,m floor (\S3.9), which these gauges
never reach. Also out of scope: the tidal classification by longitude (\S3.2), which
is a lookup; the statistics module, which is the Timeseries Statistics Model's subject; and whether synthetic
histories are a sound stand-in for observed ones (\S7.4). This is a benchmark of the
implementation (validation question Q8), not independent validation (Q9).

% =========================================================================
\mkmframework
% =========================================================================

Following the document, a day's level is
$W_t=\mu_0+S_t+T_t+\varepsilon_t$ unless a flood is injected, with
$\mu_0=h_{\mathrm{alert}}U(0.3,0.5)$, $S_t=A_s\cos(2\pi(d_t-15)/365)$,
$T_t=A_\tau M_t V_t$ with $V_t\sim U(-1,1)$, and $\varepsilon_t$ an AR(1) with
coefficient $\phi$ and marginal deviation $\sigma_w$. A severe flood occurs with daily
probability $p=\max(n_3,1)/N$ and a minor one with probability $5p$. Each sets the
level uniformly in its band, and resets $\varepsilon$ to $R=W_t-\mu_0$.

\paragraph{The exact expectation.} If the last flood was $k$ days earlier, then
$\varepsilon_t=\phi^kR+\sigma_w\sqrt{1-\phi^{2k}}\,Z$. So, given $\mu_0$, the day's
level is normal about $\mu_0+S_t+\phi^kR+T_t$, with $R$ and $V_t$ uniform. The chance
that it reaches a threshold $h$ is
\begin{equation}
P(W_t\ge h\mid k)=\mathbb{E}\Bigl[Q\Bigl(\tfrac{h-\mu_0-S_t-\phi^kR-A_\tau M_tV_t}
{\sigma_w\sqrt{1-\phi^{2k}}}\Bigr)\Bigr],
\end{equation}
a normal tail $Q$ averaged over two uniforms. Integrating twice gives a closed form in
$G(z)=\int_z^\infty Q$ and $H(z)=\int_z^\infty G$. The age $k$ is geometric, with the
flood probabilities as its parameter, and the levels are rounded to the millimetre as
reported. Summing over the run's calendar gives each expected count and the expected
mean, January and July levels. The base level $\mu_0$ is then integrated by
24-point Gauss--Legendre.

% =========================================================================
\section{Benchmark Method}
% =========================================================================

Each gauge is generated \bmReps{} times, with seeds 1 to \bmReps{} passed explicitly,
so that the run does not depend on the per-process hash. Each statistic's mean over
the series is compared with its exact expectation in standard errors, and a check
passes at $|z|\le4$. The six gauges share their seeds, and so their base-level draws,
so their errors are correlated. That is why the mean level sits about one standard
error low on all six at once. The reference is evaluated in double precision.
Wolfram derives the double integral above symbolically with \texttt{Integrate}, from
each input's exact value, and evaluates it to 30 digits at 18 points (three regimes,
six days and ages), down to tails of $10^{-53}$. There it is compared with the
reference's closed form.

The claims are judged from the exact expectations, not from the samples, except for
reproducibility, which hashes one gauge id in two fresh Python processes with
different \texttt{PYTHONHASHSEED}s.

% =========================================================================
\mkmvalidation
% =========================================================================

% Generated by tools.wolfram_oracle -- do not edit.
\begin{table}[H]\centering\small
\caption{Benchmark checks}\label{tab:bm-checks}
\begin{adjustbox}{max width=\textwidth}
\begin{tabular}{p{88mm}p{18mm}rl}
\toprule
Check & Result & Limit & Outcome \\
\midrule
Days at or above FloodAlert: model against the exact expectation, worst $|z|$ & 1.39 & 4 & pass \\
Days at or above FloodWarning: model against the exact expectation, worst $|z|$ & 2.03 & 4 & pass \\
Days at or above SevereFloodWarning: model against the exact expectation, worst $|z|$ & 1.92 & 4 & pass \\
Mean level: model against the exact expectation, worst $|z|$ & 1.02 & 4 & pass \\
January mean less the annual mean: model against the exact expectation, worst $|z|$ & 1.37 & 4 & pass \\
July mean less the annual mean: model against the exact expectation, worst $|z|$ & 1.40 & 4 & pass \\
Wolfram's exact integrals agree with the closed forms & $4.9\times10^{-10}$ & $1.0\times10^{-9}$ & pass \\
\bottomrule
\end{tabular}
\end{adjustbox}
\end{table}


\subsection{The gauges}
% Generated by tools.wolfram_oracle -- do not edit.
\begin{table}[H]\centering\small
\caption{Exact expectation / model mean over 200 series, per 50-year series}\label{tab:bm-gauges}
\begin{adjustbox}{max width=\textwidth}
\begin{tabular}{lrrrrr}
\toprule
Gauge & Floods & Days $\ge$ alert & $\ge$ warning & $\ge$ severe & Mean level (m) \\
\midrule
Non-tidal/n3=0 & 6.0 & 16.26 / 15.82 & 2.62 / 2.42 & 1.02 / 0.96 & 1.206 / 1.193 \\
Non-tidal/n3=8 & 48.0 & 129.67 / 129.42 & 20.93 / 20.68 & 8.19 / 8.14 & 1.245 / 1.232 \\
Partially tidal/n3=0 & 6.0 & 10.52 / 10.12 & 1.70 / 1.48 & 1.00 / 0.88 & 1.204 / 1.191 \\
Partially tidal/n3=8 & 48.0 & 84.04 / 84.02 & 13.55 / 13.31 & 8.00 / 7.88 & 1.227 / 1.215 \\
Tidal/n3=0 & 6.0 & 7.64 / 7.32 & 1.06 / 0.96 & 1.00 / 0.88 & 1.202 / 1.190 \\
Tidal/n3=8 & 48.0 & 61.00 / 61.09 & 8.49 / 8.24 & 8.00 / 7.88 & 1.217 / 1.204 \\
\bottomrule
\end{tabular}
\end{adjustbox}
\end{table}


\subsection{The model document's claims}
% Generated by tools.wolfram_oracle -- do not edit.
\begin{table}[H]\centering\small
\caption{The model document's claims, tested}\label{tab:bm-claims}
\begin{adjustbox}{max width=\textwidth}
\begin{tabular}{p{20mm}p{42mm}p{74mm}l}
\toprule
Claim & Stated & Benchmark finds & Outcome \\
\midrule
\S7.1(1) & lag-1 autocorrelation $\approx0.85$ (non-tidal) & 0.981 without floods; the model gives 0.967 and 0.916: the seasonal cycle dominates & \textbf{not met} \\
\S7.1(3) & days above FloodAlert $\approx 6\max(n_3,1)$ & 16.3 and 129.7 for $n_3=0, 8$, against 6 and 48 & \textbf{not met} \\
\S8.4 & single-day flood spikes & days above FloodAlert per flood: 2.70, 1.75, 1.27 (non-tidal, partial, tidal) & \textbf{not met} \\
\S5.6 & January, July means $\approx\pm A_s$ & 0.985--0.988 of $A_s$ & met \\
\S4.5 & reproduces the observed severe frequency & $n_3$ counts 5 years but is spread over 50: 8 observed per 5 years become 0.8 & \textbf{not met} \\
\S5.4, \S7.3 & byte-identical regeneration & one gauge id seeds 4181464545 in one process and 3341949591 in another; the series also ends at the run date & \textbf{not met} \\
\S6.1 & lag-7 autocorrelations 0.32, 0.028 & 0.3206, 0.0280 & met \\
\S7.2 & mean levels below FloodAlert & at most 1.245\,m & met \\
\S3.5 & a 14.8-day spring--neap cycle & its phase is the day of the year, so it jumps 9.8 days each 1 January & finding \\
\bottomrule
\end{tabular}
\end{adjustbox}
\end{table}


% =========================================================================
\mkmlimitations
% =========================================================================

\begin{enumerate}[leftmargin=*]
\item \textbf{Benchmark, not validation.} The code is shown to generate its documented
  process. Whether that process resembles a Thames gauge is \S7.4's backtest, which
  this does not perform.
\item \textbf{Severe frequency.} Eq.~(9) should divide $n_3$ by the five years it
  counts, not by $Y$; or \S4.5 and Table~2 should say what $n_3$ is meant to be. This
  raises every gauge's flood rate tenfold, and with it the upper tail of every
  generated history, so it is for the model owner and the MRC.
\item \textbf{Reproducibility.} A stable seed (for example from \texttt{hashlib}), and
  a fixed end date taken from the portfolio rather than the clock.
\item \textbf{Diagnostics.} \S7.1(1) should test the AR(1) residual, not the series.
  \S7.1(3) and \S8.4 should allow for the days a flood takes to decay. \S3.5 should
  either take the phase from a running day count or say it restarts each year.
\end{enumerate}

\mkmhistorypage{28-Oct-2026 & 0.1 & First issue & David K Kelly \\}

\end{document}
